% !TEX root = ../../main.tex
\hypearticle
{Como a IA impulsiona a originalidade da arte humana}
{Sabrina Cristan}
{

\hspace*{1.5em} Nos últimos anos, as limitações da 
inteligência 
artificial 
vêm diminuindo 
rapidamente. Hoje, a reprodução de pinturas 
e fotografias por esses modelos pode ser 
quase indistinguível da criação artística 
humana. Um trabalho manual que levaria 
horas de dedicação para ser finalizado, 
poderia ser gerado por IA em pouquíssimo 
tempo. Essa constatação leva à uma reflexão 
essencial: o que torna a arte humana algo 
verdadeiramente único? 

\articleimage
{1.0}
{sections/artigo2/arte1.png}
{Paisagem fotografada por Alvin Greis.}

\articleimage
{1.0}
{sections/artigo2/arte2.png}
{Fotografia gerada por IA.}

\subtitle{Espontaneidade como essência}

\hspace*{1.5em} A resposta parece estar na 
espontaneidade. Por mais poderosa que seja, 
a IA opera com base em previsões, e o 
imprevisível escapa ao olhar de qualquer 
algoritmo. Se uma imagem, mesmo que 
complexa, pode ser reproduzida por tal 
tecnologia, significa que a sua composição 
não é, afinal, tão original assim, pois pode ser 
reduzida a uma "fórmula".  
\hspace*{1.5em} Nesse sentido, artistas podem sentir 
desconforto ao terem a autenticidade dos 
seus trabalhos verificada por IAs, como se eles 
estivessem sendo, aos poucos, substituídos 
por um criador mais rápido e preciso. 

\subtitle{Parceria criativa}
\hspace*{1.5em} Por outra perspectiva, sem o avanço da 
inteligência artificial, o rigor artístico poderia 
não evoluir tanto. As IAs podem permitem que 
artistas se reinventem, capturam de forma 
ainda mais apurada a essência humana em 
suas obras, fujam do óbvio e criem algo 
extraordinariamente singular, desafiando, ou 
até mesmo unindo-se, às alternativas que a 
inteligência artificial traz ao mundo artístico. 

\articleimage
{1.0}
{sections/artigo2/arte3.png}
{O píer, os carros, as árvores, todos sustentam a fotografia em sua essência. Mas o gesto intencional da câmera abre uma fratura na própria realidade. Algo que os algoritmos ainda não sabem repetir.}

\hspace*{1.5em} Nesse sentido, vale mencionar a artista 
Holly Herndon que, ao lado do seu parceiro 
Mathew Dryhurst, explora as possibilidades de 
intersecção entre arte e machine learning, 
principalmente através da música, onde a 
inteligência 
artificial 
tem um papel 
fundamental como parceira criativa. Herndon 
treinou um modelo de machine learning com 
uma vasta variedade de fonemas da língua 
inglesa, com o intuito de possibilitar que o seu 
software, batizado de Holly+, pudesse 
reproduzir muitos sons diferentes. Após esse 
treinamento, a artista conseguiu mapear sua 
voz para outros idiomas e atingir tons além da 
capacidade física da voz da Holly humana. 
Atualmente, ela utiliza essa ferramenta em 
suas performances musicais.


\subtitle{Conclusão}
\hspace*{1.5em} Apesar dos avanços, ainda há muitas 
criações que não podem ser reproduzidas, 
pois não foram rotuladas nem transformadas 
em dados para alimentar um algoritmo. Além 
disso, com tantos avanços tecnológicos, as 
possibilidades da arte são reafirmadas e, de 
inimiga, a tecnologia surge como uma aliada. 
E é justamente aí que reside a verdadeira 
"beleza" da IA: uma ferramenta que incentiva 
as pessoas a se arriscarem mais e a se 
embriagarem de criatividade, capacidade 
humana que nenhum algoritmo poderia 
replicar tão bem.

\originalsource{https://medium.com/full-frame/liberation-from-the-predictable-8d65c25e3b18} 
    \textcolor{HypePurple2}{Outras fontes:}\\
\textbf{\textcolor{HypePurple2}{- }}\href{https://www.sp-arte.com/editorial/a-arte-diante-da-inteligencia-artificial}{Sp-arteSP-Arte}\\
\textbf{\textcolor{HypePurple2}{- }}\href{https://www.npr.org/transcripts/1119220726}{NprHolly Herndon: How AI can transform your voice}


}